% ============================================
% 五、编程题
% ============================================
\section{编程题：电梯调度系统}

\begin{frame}{编程题}{题目背景}
    \begin{exampleblock}{场景}
        \footnotesize 致远大厦有一部可停靠1$\sim$26层楼的电梯。\\
        设计电梯调度系统，模拟考试分值：\textbf{4+2+4=10分}
    \end{exampleblock}
    \pause
    \begin{block}{输入格式}
        \footnotesize
        \begin{itemize}
            \setlength{\itemsep}{0pt}
            \item 第一行：\texttt{n a d} —— n个按钮按下，a当前楼层，d方向（1上行/-1下行）
            \item 第二行：\texttt{n}个整数 —— 已按下的按钮（不重复、不含a、不一定有序）
        \end{itemize}
    \end{block}
    \pause
    \begin{block}{输出要求}
        \footnotesize 输出下一个应停靠的楼层。
        \begin{itemize}
            \setlength{\itemsep}{0pt}
            \item 当前方向有目标→输出该方向最近的
            \item 当前方向无目标→反方向就近停靠
            \item 无任何按钮按下→停靠1楼
        \end{itemize}
    \end{block}
\end{frame}

\begin{frame}[fragile]{编程题（一）代码}{my\_sort — 排序函数（4分）}
    \begin{exampleblock}{void my\_sort(int *arr, int cnt)}
        对数组 arr 进行\underline{原地升序}排序，元素个数为 cnt。
    \end{exampleblock}
    \begin{lstlisting}[language=C]
void my_sort(int *arr, int cnt) {
    for (int i = 0; i < cnt - 1; i++) {
        for (int j = 0; j < cnt - i - 1; j++) {
            if (arr[j] > arr[j + 1]) {
                int temp = arr[j];
                arr[j] = arr[j + 1];
                arr[j + 1] = temp;
            }
        }
    }
}
\end{lstlisting}
\end{frame}

\begin{frame}{编程题（一）解析}{my\_sort — 排序函数}
    \begin{alertblock}{解析}
        经典冒泡排序，时间复杂度 O($n^2$)。\\
        {\color{highlight}排序是为了后续能快速找到"比target大但最小的元素"。}
    \end{alertblock}
\end{frame}

\begin{frame}[fragile]{编程题（二）代码}{find\_larger — 查找函数（2分）}
    \begin{exampleblock}{int find\_larger(int *arr, int cnt, int target)}
        找到比target大、但比数组中其他元素都小的元素（即大于target的最小值）。
    \end{exampleblock}
    \begin{lstlisting}[language=C]
#include <limits.h>  // INT_MAX

int find_larger(int *arr, int cnt, int target) {
    int min_larger = INT_MAX;  // 初始化为最大整数
    for (int i = 0; i < cnt; i++) {
        if (arr[i] > target && arr[i] < min_larger)
            min_larger = arr[i];
    }
    if (min_larger == INT_MAX) return -1;  // 未找到
    return min_larger;
}
\end{lstlisting}
\end{frame}

\begin{frame}{编程题（二）解析}{find\_larger — 查找函数}
    \begin{alertblock}{解析}
        遍历数组，记录满足"大于target"条件中的最小值。\\
        find\_smaller 与之对称（找小于target的最大值），题目给出声明可直接调用。
    \end{alertblock}
\end{frame}

\begin{frame}[fragile]{编程题（三）}{main 函数 — 电梯调度逻辑（4分）}
    \begin{lstlisting}[language=C, basicstyle=\tiny\ttfamily]
int main() {
    int n, a, d;
    scanf("%d %d %d", &n, &a, &d);
    if (n == 0) { printf("1\n"); return 0; }          // 无按钮→停1楼

    int arr[100];
    for (int i = 0; i < n; i++) scanf("%d", &arr[i]);
    my_sort(arr, n);                                    // 先排序

    int next_floor;
    if (d == 1) {                                       // 上行
        next_floor = find_larger(arr, n, a);            // 找上方最近的
        if (next_floor == -1)                           // 上方无按钮
            next_floor = find_smaller(arr, n, a);       // 反向：下方最近的
    } else {                                            // 下行
        next_floor = find_smaller(arr, n, a);           // 找下方最近的
        if (next_floor == -1)                           // 下方无按钮
            next_floor = find_larger(arr, n, a);        // 反向：上方最近的
    }
    printf("%d\n", next_floor);
    return 0;
}
\end{lstlisting}
\end{frame}

\begin{frame}{编程题}{调度逻辑图解}
    \begin{block}{核心思路}
        \footnotesize
        \begin{enumerate}
            \setlength{\itemsep}{0pt}
            \item 排序按钮数组 → 有序后便于查找
            \item 上行(d=1)：找大于当前楼层的最小按钮（上方最近）
            \item 下行(d=-1)：找小于当前楼层的最大按钮（下方最近）
            \item 当前方向找不到目标 → 折返，反方向就近停靠
        \end{enumerate}
    \end{block}
    \pause
    \begin{exampleblock}{举例}
        \footnotesize
        当前在5楼，上行，按钮有 \{2, 8, 10\}：
        \begin{itemize}
            \setlength{\itemsep}{0pt}
            \item 上行方向：大于5的最小值是8 → 停8楼
        \end{itemize}

        当前在10楼，上行，按钮有 \{2, 3, 6\}：
        \begin{itemize}
            \setlength{\itemsep}{0pt}
            \item 上行方向：大于10的值不存在 → 折返
            \item 下行方向：小于10的最大值是6 → 停6楼
        \end{itemize}
    \end{exampleblock}
\end{frame}
